\documentclass[11pt,a4paper]{article}
\usepackage[utf8]{inputenc}
\usepackage[T1]{fontenc}
\usepackage[french]{babel}
\usepackage{lmodern}
\usepackage[margin=2.2cm]{geometry}
\usepackage{booktabs,array,tabularx}
\usepackage{xcolor}
\usepackage{fancyvrb}
\usepackage[hidelinks]{hyperref}
\usepackage{textcomp}

\setlength{\parskip}{0.4em}
\setlength{\parindent}{0pt}
\fvset{fontsize=\small,frame=single,framesep=2mm,rulecolor=\color{black!30}}
\newcommand{\cmd}[1]{\texttt{#1}}
\newcommand{\aremplacer}[1]{\textcolor{red!60!black}{\texttt{#1}}}

\title{Héberger l'application web \emph{VetTwins} sur un Mac mini\\[0.3em]
\large Accès réservé au réseau Tailscale (tailnet)}
\author{Projet VetTwin}
\date{17 novembre 2026}

\begin{document}
\sloppy
\maketitle

\begin{center}
\fbox{\parbox{0.92\linewidth}{\textbf{Usage pédagogique uniquement.} VetTwins est un patient virtuel canin destiné à l'enseignement. Aucun résultat ne doit servir à décider d'une dose ou d'un traitement réel.}}
\end{center}

\section{Principe}
L'application (\cmd{vetsim-kit/web/}) est un site \emph{statique} : le moteur Rust est compilé en WebAssembly et tout le calcul se fait dans le navigateur du visiteur. Le Mac mini n'a donc qu'à servir des fichiers. Le montage retenu :
\begin{enumerate}
  \item le Mac mini compile l'application (\cmd{sh web/build.sh}) ;
  \item un petit serveur HTTP (Python, lancé par \cmd{launchd} au démarrage) sert \cmd{web/} sur \cmd{127.0.0.1:8000}, invisible depuis le réseau local ;
  \item Tailscale publie ce port sur le tailnet, en HTTPS avec un certificat valide (\cmd{tailscale serve}). Seules les machines du tailnet y accèdent ; rien n'est exposé sur Internet.
\end{enumerate}
Une variante sans HTTPS, par port ouvert directement sur l'adresse Tailscale, est donnée en section~\ref{sec:portdirect}.

Dans la suite, remplacez \aremplacer{VOUS} par votre nom d'utilisateur macOS (\cmd{whoami}). Toutes les commandes se tapent dans le Terminal du Mac mini (en local ou par \cmd{ssh}).

\section{Prérequis sur le Mac mini}
\subsection{Outils de compilation}
\begin{Verbatim}
xcode-select --install                     # compilateur, git, python3
/bin/bash -c "$(curl -fsSL \
  https://raw.githubusercontent.com/Homebrew/install/HEAD/install.sh)"
curl --proto '=https' --tlsv1.2 -sSf https://sh.rustup.rs | sh -s -- -y
source "$HOME/.cargo/env"
rustup target add wasm32-unknown-unknown
cargo install wasm-bindgen-cli --version 0.2.100 --locked
\end{Verbatim}
Installez Rust par \cmd{rustup} et non par \cmd{brew install rust} : seule la version \cmd{rustup} sait ajouter la cible \cmd{wasm32-unknown-unknown}. La version de \cmd{wasm-bindgen-cli} doit être \emph{exactement} celle du \cmd{Cargo.lock} (0.2.100), sinon \cmd{build.sh} échoue avec un message de schéma incompatible.

\subsection{Mac mini toujours disponible}
\begin{Verbatim}
sudo pmset -a sleep 0 disksleep 0          # pas de mise en veille
sudo pmset -a autorestart 1                # redémarre après une coupure de courant
\end{Verbatim}

\subsection{Tailscale en service système}
Pour un serveur sans session ouverte, utilisez le démon Tailscale de Homebrew, qui démarre avec la machine (l'application Tailscale du Mac ne tourne que lorsqu'un utilisateur est connecté) :
\begin{Verbatim}
brew install tailscale
sudo brew services start tailscale
sudo tailscale up --hostname=vettwins
\end{Verbatim}
Si l'application Tailscale du Mac est déjà installée, quittez-la et désactivez son lancement à l'ouverture de session : les deux ne doivent pas tourner ensemble.

Dans la console d'administration Tailscale (\url{https://login.tailscale.com/admin}) :
\begin{itemize}
  \item \textbf{Machines} \textrightarrow{} \cmd{vettwins} \textrightarrow{} \emph{Disable key expiry}, sinon le serveur quitte le tailnet au bout de 180 jours ;
  \item \textbf{DNS} : activer \emph{MagicDNS} et \emph{HTTPS Certificates} (nécessaire à \cmd{tailscale serve} en HTTPS).
\end{itemize}

\section{Récupérer et compiler l'application}
Placez le dépôt hors de \cmd{Documents}, \cmd{Bureau} et \cmd{Téléchargements} : macOS interdit à un service \cmd{launchd} de lire ces dossiers, ce qui donne des erreurs 403 ou \og Operation not permitted \fg{} difficiles à diagnostiquer.
\begin{Verbatim}
mkdir -p ~/srv && cd ~/srv
git clone https://github.com/vinlmoe/VetTwin.git
cd VetTwin
git checkout m0-m1-workspace-dsl
cd vetsim-kit
sh web/build.sh
\end{Verbatim}
Le dépôt étant privé, \cmd{git clone} demande une authentification : connectez-vous avec \cmd{gh auth login} (\cmd{brew install gh}) ou utilisez une clé SSH (\cmd{git@github.com:vinlmoe/VetTwin.git}).

\cmd{build.sh} produit, dans \cmd{web/} : \cmd{pkg/} (moteur WebAssembly et liaisons JS), \cmd{modules.json} et \cmd{selftest.json}. Ces fichiers ne sont pas versionnés ; il faut relancer \cmd{build.sh} après chaque mise à jour. Le script n'utilise que la bibliothèque standard de Python (pas de \cmd{pip install}).

Test local rapide :
\begin{Verbatim}
python3 -m http.server --bind 127.0.0.1 -d web 8000
# autre terminal :
curl -sI http://127.0.0.1:8000/pkg/vetsim_wasm_bg.wasm | grep -i content-type
# attendu : Content-type: application/wasm
\end{Verbatim}
Arrêtez ensuite ce serveur (\cmd{Ctrl-C}) : \cmd{launchd} prendra le relais.

\section{Serveur HTTP lancé au démarrage}
Créez \cmd{/Library/LaunchDaemons/fr.vettwins.web.plist} (\cmd{sudo nano \ldots}) avec le contenu suivant, en remplaçant \aremplacer{VOUS} aux deux endroits :
\begin{Verbatim}
<?xml version="1.0" encoding="UTF-8"?>
<!DOCTYPE plist PUBLIC "-//Apple//DTD PLIST 1.0//EN"
  "http://www.apple.com/DTDs/PropertyList-1.0.dtd">
<plist version="1.0">
<dict>
  <key>Label</key>            <string>fr.vettwins.web</string>
  <key>UserName</key>         <string>VOUS</string>
  <key>ProgramArguments</key>
  <array>
    <string>/usr/bin/python3</string>
    <string>-m</string> <string>http.server</string>
    <string>--bind</string> <string>127.0.0.1</string>
    <string>--directory</string>
    <string>/Users/VOUS/srv/VetTwin/vetsim-kit/web</string>
    <string>8000</string>
  </array>
  <key>RunAtLoad</key>        <true/>
  <key>KeepAlive</key>        <true/>
  <key>StandardOutPath</key>  <string>/tmp/vettwins-web.log</string>
  <key>StandardErrorPath</key><string>/tmp/vettwins-web.log</string>
</dict>
</plist>
\end{Verbatim}
Le service tourne sous votre compte (\cmd{UserName}) et non en \cmd{root}. Activation :
\begin{Verbatim}
sudo chown root:wheel /Library/LaunchDaemons/fr.vettwins.web.plist
sudo chmod 644 /Library/LaunchDaemons/fr.vettwins.web.plist
sudo launchctl bootstrap system /Library/LaunchDaemons/fr.vettwins.web.plist
sudo launchctl print system/fr.vettwins.web | grep state     # attendu : running
curl -s -o /dev/null -w "%{http_code}\n" http://127.0.0.1:8000/   # attendu : 200
\end{Verbatim}
Arrêt définitif : \cmd{sudo launchctl bootout system/fr.vettwins.web}.

\section{Publication sur le tailnet (HTTPS, recommandé)}
\begin{Verbatim}
sudo tailscale serve --bg 8000
tailscale serve status
\end{Verbatim}
\cmd{tailscale serve status} affiche l'adresse, de la forme
\cmd{https://vettwins.\aremplacer{votre-tailnet}.ts.net/}. Le réglage est conservé après redémarrage (option \cmd{-{}-bg}). Le premier accès peut prendre quelques secondes, le temps que Tailscale obtienne le certificat.

Ce mode n'ouvre aucun port sur le réseau local et ne demande aucune règle du pare-feu macOS : Tailscale reçoit la connexion et la relaie vers \cmd{127.0.0.1:8000}.

Retrait : \cmd{sudo tailscale serve -{}-https=443 off} (ou \cmd{sudo tailscale serve reset}).

\textbf{N'utilisez pas} \cmd{tailscale funnel} : il publierait l'application sur Internet.

\section{Variante : port ouvert directement sur l'adresse Tailscale}\label{sec:portdirect}
Si vous préférez une adresse \cmd{http://vettwins:8000} sans HTTPS :
\begin{enumerate}
  \item relevez l'adresse Tailscale du Mac mini : \cmd{tailscale ip -4} (de la forme \cmd{100.x.y.z}) ;
  \item dans le plist, remplacez \cmd{127.0.0.1} par cette adresse, puis rechargez le service :
\begin{Verbatim}
sudo launchctl bootout system/fr.vettwins.web
sudo launchctl bootstrap system /Library/LaunchDaemons/fr.vettwins.web.plist
\end{Verbatim}
  \item si le pare-feu macOS est actif (\emph{Réglages Système} \textrightarrow{} \emph{Réseau} \textrightarrow{} \emph{Coupe-feu}), autorisez les connexions entrantes pour \cmd{python3}.
\end{enumerate}
Le serveur n'écoute alors que sur l'interface Tailscale : il reste invisible du Wi-Fi ou du réseau de l'école. Au démarrage, si Tailscale n'a pas encore son adresse, Python échoue puis \cmd{KeepAlive} le relance jusqu'à ce que l'adresse existe. N'utilisez pas \cmd{0.0.0.0}, qui exposerait le port sur tous les réseaux.

\section{Restreindre l'accès (facultatif)}
Par défaut, toute machine du tailnet peut ouvrir l'application. Pour limiter l'accès à un groupe, étiquetez le serveur (\cmd{sudo tailscale up -{}-hostname=vettwins -{}-advertise-tags=tag:vettwins}) et ajoutez à la politique d'accès (\emph{Access controls} dans la console) :
\begin{Verbatim}
"tagOwners": { "tag:vettwins": ["autogroup:admin"] },
"groups":    { "group:enseignants": ["prenom.nom@exemple.fr"] },
"grants": [
  { "src": ["group:enseignants"], "dst": ["tag:vettwins"],
    "ip":  ["tcp:443"] }          // "tcp:8000" pour la variante port direct
]
\end{Verbatim}
Si votre politique contient encore la règle par défaut qui autorise tout, cette restriction n'a d'effet qu'une fois cette règle retirée. Pour des étudiants hors de votre tailnet, Tailscale permet aussi de \emph{partager} la seule machine \cmd{vettwins} (\emph{Machines} \textrightarrow{} \emph{Share}).

\section{Vérification depuis un autre appareil du tailnet}
\begin{enumerate}
  \item Ouvrir \cmd{https://vettwins.\aremplacer{votre-tailnet}.ts.net/} : l'en-tête doit afficher la version de l'API du moteur, puis \og Évaluer \fg{} doit produire un patient.
  \item Ouvrir \cmd{/selftest.html} : le navigateur recalcule les cas \cmd{eval} du golden et les compare au binaire natif ; toutes les valeurs doivent être identiques au bit près.
  \item Un lien de cas partageable, par exemple :\\
  {\small\cmd{\#disease=mmvd\&stage=C\&tab=treatment\&plan=furosemide:2:po:12:14\&days=7\&sim=1}}
\end{enumerate}

\section{Mise à jour}
\begin{Verbatim}
cd ~/srv/VetTwin && git pull
cd vetsim-kit && sh web/build.sh
\end{Verbatim}
Inutile de redémarrer le service : les fichiers sont relus à chaque requête. Côté visiteur, un rechargement forcé (\cmd{Cmd-Maj-R}) évite de garder l'ancien moteur en cache.

\section{Dépannage}
\begin{tabularx}{\linewidth}{>{\raggedright\arraybackslash}p{0.36\linewidth}X}
\toprule
\textbf{Symptôme} & \textbf{Cause probable et remède} \\
\midrule
\og Chargement du moteur\ldots{} \fg{} sans fin & \cmd{web/pkg/} absent : relancer \cmd{sh web/build.sh}. Vérifier la console du navigateur. \\
\cmd{build.sh} : erreur de schéma \cmd{wasm-bindgen} & Version de \cmd{wasm-bindgen-cli} différente de 0.2.100 : réinstaller avec \cmd{-{}-version 0.2.100 -{}-locked}. \\
\cmd{build.sh} : \cmd{can't find crate for core} & Cible absente : \cmd{rustup target add wasm32-unknown-unknown}. \\
Service \cmd{launchd} arrêté, \og Operation not permitted \fg{} dans \cmd{/tmp/vettwins-web.log} & Dépôt dans un dossier protégé (Documents, Bureau) : le déplacer dans \cmd{\textasciitilde/srv}. \\
\cmd{Address already in use} & Un autre serveur occupe le port 8000 : \cmd{lsof -i :8000}. \\
\cmd{tailscale serve} refuse HTTPS & \emph{HTTPS Certificates} non activé dans la console (section DNS). \\
Le site ne répond plus après quelques mois & Expiration de la clé de la machine : \emph{Disable key expiry}, puis \cmd{sudo tailscale up}. \\
Inaccessible après redémarrage & Vérifier \cmd{sudo brew services list} (tailscale \cmd{started}) et \cmd{sudo launchctl print system/fr.vettwins.web}. \\
\bottomrule
\end{tabularx}

\end{document}
